\section{Omitted Proofs}
\label{sec:omit}


\subsubsection*{Proof of \Cref{claim:z}}
	The proof resembles the proof of Claim 3.10 in \cite{coester2022competitive}. Assume that there are $t \in T$ and $S \subseteq \cP$ such that $x'$ violates the constraint of the LP \eqref{eq:LP2} corresponding to $t$ and  $S$ and that there is $q \in \cP \setminus S$ such that $x'_q(r(q,t)) = 1$. We show that $x'$ also violates the constraint corresponding to $t$ and $S+q$.  %As $x'$ satisfies all minimal constraints, there is 
	\begin{equation*}
		\label{eq:PQPQPQ}
		\begin{aligned}
	 N-g(S) >{} & \sum_{p \in \cP-p_t} x'_p(r(p,t)) \cdot g_S(p)
	 = x'_q(r(q,t)) \cdot g_S(q)+\sum_{p \in \cP-p_t-q} x'_p(r(p,t)) \cdot g_S(p)\\
	 ={} & g_S(q)+\sum_{p \in \cP-p_t-q} x'_p(r(p,t)) \cdot g_S(p)
	  \geq g_S(q)+\sum_{p \in \cP-p_t-q} x'_p(r(p,t)) \cdot g_{S+q}(p).\\
		\end{aligned}
	\end{equation*} The first inequality holds since we assume that $x'$ violates the primal constraint corresponding to $t$ and $S$. The second equality follows because $x'_q(r(q,t)) = 1$. The second inequality follows from the submodularity of $g$. As $g(S+q) = g(S)+g_{S}(q)$, by the above inequality we conclude that $x'$ also violates the constraint corresponding to $t$ and $S+q$. Hence, if $x'$ satisfies all minimal constraints it also satisfies all non-minimal constraints. \qed



\subsubsection*{Proof of \Cref{lem:primalX}}
 By \Cref{claim:z} we only need to show that $x$ satisfies all minimal constraints. %Observe that \Cref{step:Whilefeasible} along with the submodularity of $g$ ensures that the algorithm returns a feasible (fractional) solution $x$ to~\eqref{eq:LP2}. Specifically, 
If $x$ is not feasible at time $t$, the algorithm is guaranteed to find a minimal set $Q$ in \Cref{step:fQ} such that the primal constraint for $t$ and $Q$ is not satisfied. Then, the algorithm increases the variables of pages in $(\cP-p_t) \setminus Q$ until they satisfy the constraint of $Q,t$ by  \Cref{step:Whilefeasible} or until $Q$ is no longer minimal; as we only need to prove that $x$ satisfies all minimal constraints by \Cref{claim:z}, either the constraint corresponding to $t,Q$ is satisfied, or this constraint is not minimal anymore. %(i.e., $Q$ is not minimal). 
Thus, $x$ satisfies all minimal constraints, yielding the proof. \qed







\subsubsection*{Proof of \Cref{lem:zFeasible}}
 	By \Cref{claim:z} we only need to show that $z$ satisfies all minimal constraints of \eqref{eq:LP2}. Let $t \in T$ and let $S \subseteq \cP$ be such that $t$ and $S$ form a minimal constraint for $z$ in  \eqref{eq:LP2}. If $g(S) = N$ then the constraint of $t$ and $S$ is trivially satisfied by $z$; thus, assume that $N-g(S) \geq 1$. As the constraint of $t$ and $S$ is minimal for $z$, for all $p \in \cP\setminus S$ it holds that $z_p(r(p,t)) < 1$; thus, by \Cref{step:1/2} it follows that $z_p(r(p,t)) \leq \frac{1}{2}$. Hence,  by \Cref{step:z} for all $p \in \cP\setminus S$ it holds that 
 $z_p(r(p,t)) \geq 2 \cdot x_p(r(p,t))-\frac{1}{4 \cdot N \cdot \mu}$.
 Therefore,
\begin{equation}
	\label{eq:z2}
	\begin{aligned}
		 \sum_{p \in \cP-p_t} g_S(p) \cdot z_p(r(p,t))
		\geq{} &   \sum_{p \in \cP-p_t} g_S(p) \cdot \left( 2 \cdot x_p(r(p,t))-\frac{1}{4 \cdot N \cdot \mu} \right)\\
		={} & 2 \cdot \sum_{p \in \cP-p_t} g_S(p) \cdot x_p(r(p,t))-\sum_{p \in \cP-p_t} g_S(p) \cdot \frac{1}{4 \cdot N \cdot \mu} \\
		\geq{} & 2 \cdot \sum_{p \in \cP-p_t} g_S(p) \cdot x_p(r(p,t))-N \cdot \mu \cdot \frac{1}{4 \cdot N \cdot \mu} \\
			\geq{} & 2 \cdot \left(N-g(S)\right)-\frac{1}{4} 
				\geq N-g(S). 
	\end{aligned}
\end{equation}
%The first inequality holds by \eqref{eq:A1}. 
The second inequality holds since the marginal contribution of any page $p$ to $S$ is bounded by the total cover demand $N$, i.e., $g_S(p) \leq N-g(S) \leq N$. The third inequality holds since $x$ is a feasible solution for the LP by \Cref{lem:primalX}.   The last inequality follows from the assumption $N-g(S) \geq 1$. By \eqref{eq:z2} the proof follows. \qed


\subsubsection*{Proof of \Cref{lem:FFeasible}}
To prove that the dual $y$ is also feasible, consider some $p \in \cP$ and $j \in [n_p]$. By the update rate of $x_p(j)$ in \Cref{step:fx} it follows that 
		\begin{equation}
		\label{eq:FromSAP}
		\frac{1}{\mu} \cdot \left(     \exp \left(  \frac{\ln (\mu+1)}{c(p)} \cdot Y_p(j)\right)-1\right) = x_p(j) \leq 1.
	\end{equation} 
	The inequality holds since $Y_p(j)$ does not increase in time $t$ if $p$ is already fully evicted at this time by \Cref{step:InWhile}; that is, once $x_p(r(p,t)) = 1$ it holds that $Q$ is no longer minimal and the algorithm chooses a new minimal set $Q'$ that includes $p$. Thus, by simplifying the expression in \eqref{eq:FromSAP} it holds that 
	$$\exp \left( \frac{\ln (\mu+1)}{c(p)} \cdot Y_p(j) \right) \leq \mu+1.$$
	by taking logarithms from both sides $Y_p(j) \leq c(p)$. Hence, $y$ is a feasible dual solution. \qed


\subsubsection*{Proof of \Cref{lem:ALGf}}
The feasibility of $z$ follows from \Cref{lem:zFeasible}. Moreover, the first property follows from \Cref{step:1/2,step:z} of the algorithm. For the second property, since $x$ and $y$ are feasible primal and dual solutions by \Cref{lem:primalX} and \Cref{lem:FFeasible}, it follows that the cost of $x$ is bounded by $O(\log (\mu))$ times the cost of $\OPT$. Thus, the proof follows from Observation~\ref{obs:z}. \qed

\subsubsection*{Proof of \Cref{lem:ALG3}}
The feasibility of the algorithm follows from \Cref{step:whileFe}. To bound the expected cost of the algorithm, consider some time step $t \in T$. We bound the expected cost paid at time $t$ by the algorithm on evictions, w.r.t. the cost paid by $z'$ on fetching. Let $z'_p(t-1)$ be the value of the variable $z'_{p_t}(r(p_t,t))$ at time $t-1$, i.e., indicating the portion of $p$ missing from cache at time $t-1$. If $z'_p(t-1) = 0$, i.e., $p_t$ is in the cache when it is requested at time $t$, then our integral solution obtained in \Cref{alg:randomizedRounding} also does not evict $p_t$. It follows that we do not incur a cost at time $t$ if $z'_p(t-1)) = 0$. Otherwise, by \Cref{lem:ALGf} it holds that $z'_p(t-1) \geq \frac{1}{4 \cdot N \cdot \mu}$. 
	
	Let $c_{\min}, c_{\max}$ be the minimum and maximum costs of pages, thus $C = \frac{c_{\max}}{c_{\min}}$.  As we assume that $z'_p(t-1) \geq \frac{1}{4 \cdot N \cdot \mu}$, it follows that the cost of the solution $z'$ increases by at least $\frac{c_{\min}}{4 \cdot N \cdot \mu}$ at time $t$ by fetching the fraction of $p_t$ missing from cache. At the same time, observe that each time that the algorithm enters \Cref{step:reach}, the total cost it pays is bounded by $c_{\max} \cdot \mu$ as it can evict at most $\mu$ pages from cache as the cache has been feasible until this time.  Moreover, by \Cref{lem:Gupta}, the probability that the algorithm enters \Cref{step:reach} in time $t$ is bounded by $\frac{1}{2 \cdot C \cdot N \cdot \mu^2}$. Thus, the expected cost paid by the algorithm at time $t$ is bounded by $\frac{c_{\max} \cdot n}{2 \cdot C \cdot N \cdot \mu^2} \leq \frac{c_{\min}}{2 \cdot N \cdot \mu}$. Therefore, at each time step the expected cost of the algorithm is bounded by a constant factor of the expected cost of $z'$. In addition, the cost of $z'$ is bounded by $\alpha$ times the cost of $z$. Therefore, the proof follows from the competitive ratio of $z$ described in \Cref{lem:ALGf}. \qed
	
	
	\subsubsection*{Proof of \Cref{thm:approximation}}
	
	using the ``round-or-separate''
	approach of \cite{gupta2020online} we can solve our LP \eqref{eq:LP2} offline approximately in polynomial time with $O(\log f(\cP))$-approximation ratio. Using \Cref{lem:Gupta}  and the analysis in \Cref{lem:ALG3}, this implies the statement of \Cref{thm:approximation}.  \qed